\newpage

\section{Proof of Theorem~\ref{thm:extreme_dominance}}
\label{appendix:proof_extreme_dominance}

Let
\[
X_{ij}=\mu_j+Z_{ij},
\]
where $Z_{ij}$ is mean-zero $\sigma^2$-sub-Gaussian. For $t<|\mu_j|$, reverse triangle inequality gives
\[
|X_{ij}|\ge |\mu_j|-|Z_{ij}|.
\]
Hence
\[
\{|Z_{ij}|<|\mu_j|-t\}\subseteq \{|X_{ij}|>t\},
\]
so
\[
\mathbb{P}(|X_{ij}|>t)\ge \mathbb{P}(|Z_{ij}|<|\mu_j|-t).
\]
Using the sub-Gaussian tail bound
\[
\mathbb{P}(|Z_{ij}|\ge u)\le 2\exp\!\left(-\frac{u^2}{2\sigma^2}\right),
\]
with $u=|\mu_j|-t$ yields
\[
\mathbb{P}(|X_{ij}|>t)
\ge
1-2\exp\!\left(-\frac{(|\mu_j|-t)^2}{2\sigma^2}\right).
\]
This proves Theorem~\ref{thm:extreme_dominance}.

\newpage

\section{Proof of Theorem~\ref{thm:dense_extreme}}

Let
\[
X_{ij} = \mu_j + Z_{ij},
\]
where $Z_{ij}$ are independent mean-zero sub-Gaussian random variables with variance proxy $\sigma^2$.

Define the exceedance indicator
\[
B_i = \mathbf{1}\{|X_{ij}| > t\},
\]
and the total exceedance count
\[
C_j(t) = \sum_{i=1}^{N} B_i.
\]

\subsection{Mean-dominated regime}

Assume $|\mu_j| > t$.

By the reverse triangle inequality
\[
|X_{ij}| = |\mu_j + Z_{ij}| \ge |\mu_j| - |Z_{ij}|.
\]

Thus if
\[
|Z_{ij}| < |\mu_j| - t,
\]
then necessarily
\[
|X_{ij}| > t.
\]

Therefore
\[
\mathbb{P}(|X_{ij}| > t)
\ge
\mathbb{P}(|Z_{ij}| < |\mu_j| - t).
\]

Because $Z_{ij}$ is sub-Gaussian,
\[
\mathbb{P}(|Z_{ij}| \ge u)
\le
2\exp\!\left(-\frac{u^2}{2\sigma^2}\right).
\]

Setting $u = |\mu_j| - t$ gives
\[
\mathbb{P}(|X_{ij}| > t)
\ge
1 - 2\exp\!\left(-\frac{(|\mu_j|-t)^2}{2\sigma^2}\right).
\]

Since the indicators $B_i$ are independent,

\[
\mathbb{E}[C_j(t)]
=
\sum_{i=1}^{N} \mathbb{E}[B_i]
=
\sum_{i=1}^{N} \mathbb{P}(|X_{ij}| > t).
\]

Using the lower bound above yields

\[
\mathbb{E}[C_j(t)]
\ge
N\left(1 - 2\exp\!\left(-\frac{(|\mu_j|-t)^2}{2\sigma^2}\right)\right).
\]

\subsection{Variance-only regime}

If $\mu_j = 0$, then $X_{ij} = Z_{ij}$ and

\[
\mathbb{P}(|X_{ij}| > t)
=
\mathbb{P}(|Z_{ij}| > t).
\]

Using the sub-Gaussian tail bound

\[
\mathbb{P}(|Z_{ij}| > t)
\le
2\exp\!\left(-\frac{t^2}{2\sigma^2}\right),
\]

we obtain

\[
\mathbb{E}[C_j(t)]
=
\sum_{i=1}^{N} \mathbb{P}(|Z_{ij}| > t)
\le
2N\exp\!\left(-\frac{t^2}{2\sigma^2}\right).
\]

\subsection{Interpretation}

The two regimes differ fundamentally:

\[
\text{Mean-dominated: } C_j(t) = \Theta(N)
\]

\[
\text{Variance-only: } C_j(t) = \mathcal{O}\!\left(N e^{-t^2/(2\sigma^2)}\right).
\]

Thus a coherent mean shift produces a dense population of extreme activations across tokens, whereas purely variance-driven extremes are exponentially rare.

This scaling difference explains why mean bias becomes the dominant driver of activation outliers in large models.